%%%%%%%%%%  scan idx 245  |  folio 234  %%%%%%%%%%

2.1.13. Sea $\ell$ un número primo. Para todo grupo conmutativo $E$,
denotemos por $T_\ell(E)$ el sistema proyectivo
\begin{equation}\tag{2.1.13.1}
 T_\ell(E)=({}_{\ell^\nu}E)_{\nu\geq0},
\end{equation}
que depende manifiestamente de manera funtorial de $E$. Las
compatibilidades de 2.1.12, para $n,n'$ de la forma
$\ell^\nu,\ell^{\nu'}$ con $\nu'\geq\nu$, muestran que los
homomorfismos $\alpha^{\ell^\nu}_{P,Q}$, al variar $\nu$, definen un
homomorfismo de sistemas proyectivos, funtorial en $P,Q$:
\begin{equation}\tag{2.1.13.2}
 \alpha^{(\ell)}_{P,Q}:
 T_\ell(P)\otimes T_\ell(Q)
 \longrightarrow T_\ell\bigl(\uTor_1(P,Q)\bigr).
\end{equation}
Conocer estos homomorfismos para todo primo $\ell$ equivale
manifiestamente a conocer los $\alpha^n_{P,Q}$ para todo entero $n$.
Las relaciones de antisimetría de 2.1.10 se convierten entonces en las
relaciones de antisimetría $\ell$-ádicas análogas, para todo $\ell$,
expresadas por la anticonmutatividad de los diagramas
\begin{equation}\tag{2.1.13.3}
\begin{tikzcd}[column sep=huge,row sep=huge]
 T_\ell(P)\otimes T_\ell(Q)
   \arrow[r]
   \arrow[d,"{\alpha^{(\ell)}_{P,Q}}"']
   \arrow[rd,phantom,"-" description]&
 T_\ell(Q)\otimes T_\ell(P)
   \arrow[d,"{\alpha^{(\ell)}_{Q,P}}"]\\
 T_\ell\bigl(\uTor_1(P,Q)\bigr)
   \arrow[r]&
 T_\ell\bigl(\uTor_1(Q,P)\bigr).
\end{tikzcd}
\end{equation}
Aquí las flechas horizontales son las flechas de simetría.

\medskip
\noindent\underline{Observación 2.1.14.}
La construcción del homomorfismo (2.1.1) y los resultados 2.1.10 y
2.1.11 que le conciernen pueden generalizarse a la situación en que
se dan un anillo conmutativo $\underline A$ sobre $T$ y dos ideales
$\underline a$ y $\underline b$ de $\underline A$.

%%%%%%%%%%  scan idx 246  |  folio 235  %%%%%%%%%%

En el caso tratado anteriormente, ambos ideales son el ideal
$n\mathbb Z_T$ del anillo constante $\mathbb Z_T$ sobre $T$. Se
define entonces un homomorfismo, funtorial en los $\underline A$-módulos
$P,Q$,
\begin{equation}\tag{2.1.14.1}
 {}_{\underline a}P\otimes{}_{\underline b}Q
 \xrightarrow{\ \alpha\ }
 \uHom_{\underline A}
 \left(
   \frac{\underline a\cap\underline b}
        {\underline a\mathbin{\cdot}\underline b},
   \uTor^{\underline A}_1(P,Q)
 \right),
\end{equation}
como compuesto de los homomorfismos
\[
\begin{split}
 {}_{\underline a}P\otimes{}_{\underline b}Q
 &\longrightarrow
 \uHom(\underline A/\underline a,P)
 \otimes
 \uHom(\underline A/\underline b,Q)\\
 &\longrightarrow
 \uHom\!\left(
   \uTor_1(\underline A/\underline a,\underline A/\underline b),
   \uTor_1(P,Q)
 \right),
\end{split}
\]
utilizando el conocido isomorfismo canónico
\begin{equation}\tag{2.1.14.2}
 \uTor_1(\underline A/\underline a,\underline A/\underline b)
 \simeq
 \frac{\underline a\cap\underline b}
      {\underline a\mathbin{\cdot}\underline b}.
\end{equation}
En realidad, el signo de 2.1.10 está ``oculto'' en la elección de este
isomorfismo, según se trabaje con la resolución
$0\longrightarrow\underline a\longrightarrow\underline A
\longrightarrow0$ de $\underline A/\underline a$, o con la
resolución $0\longrightarrow\underline b\longrightarrow\underline A
\longrightarrow0$ de $\underline A/\underline b$. Aquí se convierte
en la anticonmutatividad del diagrama
\begin{equation}\tag{2.1.14.3}
\begin{tikzcd}[column sep=huge,row sep=huge]
 {}_{\underline a}P\otimes{}_{\underline b}Q
   \arrow[r]
   \arrow[d]
   \arrow[rd,phantom,"-1" description]&
 {}_{\underline b}Q\otimes{}_{\underline a}P
   \arrow[d]\\
 \uHom\!\left(
   \dfrac{\underline a\cap\underline b}
         {\underline a\mathbin{\cdot}\underline b},
   \uTor_1(P,Q)
 \right)
   \arrow[r]&
 \uHom\!\left(
   \dfrac{\underline a\cap\underline b}
         {\underline a\mathbin{\cdot}\underline b},
   \uTor_1(Q,P)
 \right).
\end{tikzcd}
\end{equation}
Las flechas horizontales son las flechas de simetría. Cuando
$\underline a=\underline b$, esta anticonmutatividad también puede
expresarse diciendo que el automorfismo de simetría en
$\uTor_1(\underline A/\underline a,\underline A/\underline a)$ es
$-\id$.

En el caso particular en que
$\underline a=\underline b=n\underline A$, con $n$ una sección de
$\underline A$, la multiplicación por $n$ induce un epimorfismo
\[
 \underline A/\underline a
 \longrightarrow
 \frac{\underline a\cap\underline a}
      {\underline a\mathbin{\cdot}\underline a}
 =\underline a/\underline a^2,
\]
y, por tanto, un \sisi{isomorfismo inyectivo}{homomorfismo inyectivo},
funtorial en el $\underline A$-módulo $T$,

%%%%%%%%%%  scan idx 247  |  folio 236  %%%%%%%%%%

\[
 \uHom(\underline a/\underline a^2,T)
 \longrightarrow{}_{\underline a}T.
\]
Así, en este caso, el homomorfismo (2.1.14.1) puede interpretarse como
un homomorfismo
\begin{equation}\tag{2.1.14.4}
 {}_nP\otimes{}_nQ
 \longrightarrow{}_n\uTor_1(P,Q),
\end{equation}
de forma más próxima a (2.1.1).

Sigamos suponiendo que $\underline a=\underline b=n\underline A$ y
que $n$ es una sección regular de $\underline A$. La sucesión
espectral (2.1.4) conserva su sentido en la situación presente y da la
siguiente sucesión exacta de términos de grado bajo, que generaliza
(2.1.5), donde $Q$ es ahora un $\underline A$-módulo \underline{anulado por}
$\underline a$:
\begin{equation}\tag{2.1.14.5}
\begin{split}
 \uTor^{\underline A}_2(P,Q)
 &\longrightarrow
 \uTor^{\underline A/n\underline A}_2(P_n,Q)
 \longrightarrow{}_nP\otimes Q
 \xrightarrow{\ \alpha\ }
 \uTor^{\underline A}_1(P,Q)\\
 &\longrightarrow
 \uTor^{\underline A/n\underline A}_1(P_n,Q)
 \longrightarrow0.
\end{split}
\end{equation}

Volvamos ahora a $\underline a$ y $\underline b$ arbitrarios, y sean
$p$ y $q$ dos secciones de $\underline A$. Obtenemos entonces el
diagrama conmutativo siguiente; (2.1.11.1) es esencialmente el caso
particular que se obtiene tomando $q=1$:
{\footnotesize
\begin{equation}\tag{2.1.14.6}
\begin{tikzcd}[row sep=large,column sep=tiny]
 &{}_{p\underline a}P\otimes{}_{q\underline b}Q
   \arrow[dl,"p\otimes q"']
   \arrow[dr,"i\otimes j"]&\\
 {}_{\underline a}P\otimes{}_{\underline b}Q
   \arrow[d,"\alpha"']&&
 {}_{pq\underline a}P\otimes{}_{pq\underline b}Q
   \arrow[d,"\alpha"]\\
 \uHom\!\left(
   \dfrac{\underline a\cap\underline b}
         {\underline a\mathbin{\cdot}\underline b},
   \uTor_1(P,Q)
 \right)
   \arrow[dr,hook,"\mathrm{incl.}"']&&
 \uHom\!\left(
   \dfrac{pq\underline a\cap pq\underline b}
         {pq\underline a\mathbin{\cdot}pq\underline b},
   \uTor_1(P,Q)
 \right)
   \arrow[dl,hook,"(pq)^*"]\\
 &\uHom\!\left(
   \dfrac{\underline a\cap\underline b}
         {pq\underline a\mathbin{\cdot}\underline b},
   \uTor_1(P,Q)
 \right)&
\end{tikzcd}
\end{equation}}

%%%%%%%%%%  scan idx 248  |  folio 237  %%%%%%%%%%

Aquí
\[
 i:{}_{p\underline a}P\longrightarrow{}_{pq\underline a}P,
 \qquad
 j:{}_{q\underline b}Q\longrightarrow{}_{pq\underline b}Q
\]
son las inclusiones canónicas, y $(pq)^*$ denota el homomorfismo
inyectivo obtenido, pasando a $\uHom$, a partir del homomorfismo
suprayectivo
\[
 \frac{\underline a\cap\underline b}
      {pq\underline a\mathbin{\cdot}\underline b}
 \longrightarrow
 \frac{pq\underline a\cap\sisi{\underline b}{pq\underline b}}
      {pq\underline a\mathbin{\cdot}pq\underline b}
\]
inducido por la multiplicación por $pq$. Cuando
$\underline a=\underline b=n\underline A$, este diagrama puede
reescribirse en una forma más sencilla, dada en (2.1.11.1) para
$q=1$.
